\subsection{Matriz de evaluación ponderada}

Con base en los criterios anteriores, se aplicó una matriz de decisión ponderada en una
escala discreta de 1 a 5, donde 5 representa el mayor grado de adecuación al objetivo
del trabajo y 1 el menor. Los puntajes corresponden a una valoración relativa respecto a
la posibilidad de integrar, instrumentar y evaluar una estrategia de compresión sin
pérdida sobre el vector de estado.

\begin{table}[H]
\centering
\caption{Evaluación de alternativas (escala 1--5) y total ponderado}
\label{tab:evaluacion_sim}
\begin{tabular}{lcccccc}
\hline
\multicolumn{1}{c}{\textbf{Criterio}} &
\textbf{Peso} &
\textbf{TMFQS} &
\textbf{QuEST} &
\textbf{Intel-QS} &
\textbf{qsim} &
\textbf{Quantum++} \\
\hline
C1 & 0.15 & 5 & 5 & 5 & 5 & 4 \\
C2 & 0.25 & 5 & 3 & 2 & 2 & 3 \\
C3 & 0.20 & 5 & 3 & 2 & 2 & 3 \\
C4 & 0.10 & 3 & 5 & 5 & 5 & 4 \\
C5 & 0.10 & 4 & 3 & 2 & 2 & 5 \\
C6 & 0.10 & 4 & 5 & 5 & 4 & 3 \\
C7 & 0.10 & 5 & 3 & 3 & 2 & 2 \\
\hline
\textbf{Total} & \textbf{1.00} &
\textbf{4.70} & 3.70 & 3.20 & 2.90 & 3.30 \\
\hline
\end{tabular}
\end{table}

La matriz muestra que las alternativas mejor orientadas a rendimiento y paralelismo no
coinciden necesariamente con las mejor posicionadas para intervenir e instrumentar el
almacenamiento del estado. En este marco, TMFQSfullstate obtiene la mayor
puntuación global, principalmente por su ventaja en los criterios de intervención sobre la
representación del estado, capacidad de instrumentación experimental y afinidad con el
estudio de estrategias de memoria.
